\section{Even components in every finite quotient}
\label{sec:selected-all-stages}

This section defines the component quotient directly in the
nilpotent group, so that the construction applies at every
nilpotency class. Continue with the canonical groups
\(Q_n=H/\mathcal F_{n+1}\) of Theorem~\ref{thm:canonical-all-stages},
for \(n\geq2\), and put
\[
 B_n=\operatorname{im}(\ker(H\longrightarrow W)\longrightarrow Q_n),
 \qquad D_n=\operatorname{im}(\mathcal F_{n+1}\longrightarrow U).
\]
A tame complement gives an action of \(\Delta\) on \(Q_n\).
For an exponent \(e\), write
\[
 \mathcal E_{n,e}=
 \{g\in B_n:\delta g\delta^{-1}=g^{\omega(\delta)^e}
                         \text{ for every }\delta\in\Delta\}.
\]
The power is defined in the procyclic subgroup generated by \(g\).
For a set \(\Sigma\) of even exponents, define
\[
 J_{n,\Sigma}=
 \overline{\left\langle\!\left\langle
       \bigcup_{\substack{e\ \mathrm{even}\\e\notin\Sigma}}
                   \mathcal E_{n,e}
                   \right\rangle\!\right\rangle}_{Q_n},
 \qquad Q_{n,\Sigma}=Q_n/J_{n,\Sigma}.
\]
Every element of each displayed set is used in this definition.

\begin{lemma}\label{lem:tame-eigen-element-lift}
Let a finite group \(\Delta\) of order prime to \(p\) act
continuously on a nilpotent pro-\(p\) group \(B\), and let
\(B\twoheadrightarrow V\) be an equivariant homomorphism to
an abelian pro-\(p\) group. For a character
\(\chi:\Delta\to\mathbb Z_p^\times\) and
\(v\in V\) satisfying \(\delta v=\chi(\delta)v\),
there is \(g\in B\) above \(v\) satisfying
\(\delta g=g^{\chi(\delta)}\) for every \(\delta\).
\end{lemma}
\begin{proof}
Lift \(v\) to \(B^{\rm ab}\) and apply the integral projector
\[
 \epsilon_\chi=\frac1{|\Delta|}
               \sum_{\delta\in\Delta}\chi(\delta)^{-1}\delta.
\]
This gives an eigenclass above \(v\). We lift it successively
through the finite lower central series of \(B\).
At one central step, let \(C\) be the central abelian kernel,
and choose a lift \(u\) of the eigen-element already obtained.
Write the central errors additively as
\[
 \delta u\delta^{-1}=k_\delta u^{\chi(\delta)},
 \qquad k_\delta\in C.
\]
Put \(b_\delta=\chi(\delta)^{-1}k_\delta\) and give \(C\)
the action \(\delta\star c=\chi(\delta)^{-1}\delta(c)\).
Multiplication of two automorphisms gives
\[
 b_{\delta\eta}=b_\delta+\delta\star b_\eta.
\]
The finite average \(c=|\Delta|^{-1}\sum_\eta b_\eta\)
therefore satisfies \(b_\delta=c-\delta\star c\).
Since \(c\) is central, replacing \(u\) by \(cu\) changes
its error to
\[
 \delta(c)+k_\delta-\chi(\delta)c=0.
\]
Thus the corrected lift is an eigen-element at this step.
Division by \(|\Delta|\) is an automorphism of every abelian
pro-\(p\) group, including its torsion subgroup. The procedure
is consequently integral at each step. The lower central
series has finite length, so it ends with the required
\(g\in B\).
\end{proof}


\begin{proof}[Proof of Theorem~\ref{thm:selected-all-stages}]
Both tame complements lie in the inverse image of the finite subgroup \(\Delta\subset\Delta\times\Gamma\), whose kernel is \(Q_n\). Profinite Schur--Zassenhaus conjugates them by an element of \(Q_n\). Conjugation sends each full eigen-element set to the corresponding set for the other complement, proving independence of its closed normal closure. For an arithmetic element \(g\), choose \(h\in Q_n\) such that \(h^{-1}g\) normalizes the complement. The quotient \(\Delta\times\Gamma\) is abelian and the complement maps injectively to it, so \(h^{-1}g\) centralizes the complement. It preserves \(B_n\) and every eigen-element set with its character label. Normality handles the factor \(h\), proving full arithmetic stability. The map \(Q_n\to Q_{n-1}\) sends a
\(\Delta\)-eigen-element to a \(\Delta\)-eigen-element of
the same character. It induces a surjection
\(Q_{n,\Sigma}\to Q_{n-1,\Sigma}\).

Let \(V_n=\operatorname{im}(\mathcal F_{n+1}\to M)\). Then
\(Q_n^{\rm ab}=M/V_n\), and the image of \(B_n\) is the
image of \(D\oplus T\). Every bad even eigen-element maps
into its indicated bad even character. Conversely, apply
Lemma~\ref{lem:tame-eigen-element-lift} to the surjection
from \(B_n\) to this abelian image. Every class in a bad
even character has such a lift. Hence the abelian image of
\(J_{n,\Sigma}\) is exactly the sum of the bad even
components of \(T/(V_n)^+\), and contains no odd class.
Since \(\mathcal F_{n+1}\subseteq \mathcal F_3\) for \(n\geq2\),
\[
                   (V_n)^+\subseteq\pi_T(Z).
\]
The proved single-component cup theorem gives
\(Z^{[e]}=0\) for \(e\in\Sigma\). Therefore the selected
even part is exactly \(T_\Sigma\), and the odd part is
\(U/D_n\), proving the displayed abelianization.
Kernel containment follows from the canonical map through
\(Q_n\), and arithmetic stability gives the Galois assertion.

For exact class, choose distinct positive odd weights
\(m,k\geq3\) congruent to \(-1\) modulo \(p-1\).
Project the full free rational outer Lie algebra to the
free Lie algebra on its two letters of weights \(m,k\),
and truncate at generator length \(n\). This gives a
homomorphism from \(Q_n\) to its nilpotent exponential group.
Both letters have tame character \(-1\). Every even Lie
word of length at most \(n\) consequently has one of the
characters in the displayed condition. An element of
\(\mathcal E_{n,e}\) with bad even \(e\) has logarithm in
a missing eigenspace, so its image is zero. The homomorphism
therefore factors through \(Q_{n,\Sigma}\).

Choose \(w\in W^{[-1]}\) for which both linear moments
are nonzero. Such a choice follows from the explicit regular
\(\Lambda\)-coordinate of this character, after clearing
one scalar denominator. The two linear images of \(w\)
and \(\gamma w\) have determinant
\[
             \kappa_m(w)\kappa_k(w)(u^k-u^m)\ne0.
\]
Choose lifts to \(H\). The iterated commutator with initial
word \(\operatorname{ad}(w)^{n-2}[w,\gamma w]\) maps to a
nonzero length-\(n\) Hall word on two independent vectors.
The quotient thus has class at least \(n\); it has class
at most \(n\) as a quotient of \(Q_n\).

For \(n\leq5\), the group \(B_n\) has class at most two.
Its integral BCH logarithm identifies its group
\(\Delta\)-eigen-elements exactly with the Lie
\(\Delta\)-eigenspaces. The defining eigen-element sets are therefore the exponential images of exactly the earlier eigenspaces. Their closed normal closures are equal. This proves the
last identification through the specified maps.
\end{proof}

In particular, the thirty characters
\(\Sigma_{30}=\{-2m:1\leq m\leq30\}\) contain all the
characters required for every \(2\leq n\leq61\).
The exact class assertion gives the strict tower
\[
 L_{2,\Sigma_{30}}\subsetneq L_{3,\Sigma_{30}}
 \subsetneq\cdots\subsetneq L_{61,\Sigma_{30}}
 \subseteq\Omega_p^*.
\]
The arithmetic premise for this application is the posted
\(K\)-group vanishing statement, with the dependency analysis
of Section~\ref{sec:dependencies}. The component theorem
itself was proved for an arbitrary \(\Sigma\) using only
finiteness of its selected Iwasawa components.
